\section{Lexical Analysis} 

  When we have code in a \texttt{.py} file and run it, the \textbf{lexical analyzer} generates a stream of \textbf{tokens} to be inputted into a parser. 

  \begin{theorem}[UTF-8 in the Lexical Analyzer]
    \label{thm:utf-8-in-the-lexical-analyzer}
    All UTF-8 characters can be parsed by the lexical analyzer. 
  \end{theorem}

  \begin{question}[Characters Outside the Source Character Set]
    Which characters aren't? 
  \end{question}

  \begin{definition}[Logical and Physical Lines]
    \label{def:logical-and-physical-lines}
 
    There are two types of lines in Python. 
    \begin{enumerate}
      \item A \textbf{logical line} is represented by the token \texttt{NEWLINE}. 
      \item A \textbf{physical line} is a sequence of characters terminated by an end-of-line (EOL) sequence. 
    \end{enumerate} 
    \begin{lstlisting}
      x = [1, 2, 3, 4]  # one logical line on one physical line 
      x = [1, 2,        # one logical line on two physical lines
           3, 4]
    \end{lstlisting}
  \end{definition} 
